% Einheit 5 von 5 – Anwendungen (bank/lineare-funktionen/e5.jsonl, zusammenbau v0.1)
%% TODO Zweigzeile Teil 1: was hier gelernt wird (Satz in Schülersprache aus dem Einheitstitel) – Einheit 5
\einheitenkopf[e5]{Einheit 5 von 5 · Anwendungen}
\zweigzeile{ab Kl. 8 · P10 oft}
\setcounter{aufgabe}{27}

%% TODO Titel: Ich-kann-Satz fehlt in der Bank (Platzhalter: Kettenname) – Nr. 28
%% TODO Anweisung: ein Satz über der ersten Teilaufgabe fehlt in der Bank – Nr. 28
\begin{aufgabe}{Anwendung}
%% TODO \rechenplatz{n}: Schrittzahl des Grundfalls fehlt in der Bank (nur bei fester Schreibform, 2.3 b)
\begin{teile}
\teil Markiere: Unterstreiche den Anfangswert, kreise die Änderung je Einheit ein. Ein Handwerker berechnet 45 € für die Anfahrt und 38 € je Arbeitsstunde.
\teil Welche Gleichung passt? Ein Fahrradverleih nimmt 8 € Grundgebühr und 3 € je Stunde. y ist der Preis in €, x die Anzahl der Einheiten. Kreuze an. \\ \kreuz{$y = 8x + 3$} \\ \kreuz{$y = 3x + 8$} \\ \kreuz{$y = 11x$} \\ \kreuz{$y = 3x - 8$}
\teil Stelle eine Gleichung auf: Ein Kino-Abo kostet einmalig 10 € und 6 € je Film. y sind die Kosten in €, x die Anzahl der Filme. y = \leerfeld
\teil Wie viel nach 12 Wochen? Tom hat 85 € gespart. Jede Woche kommen 15 € dazu. \leerfeld[€]
\teil Welcher Tarif ist für 60 km günstiger? Tarif A: 20 € Grundgebühr und 0{,}30 € je km. Tarif B: 0{,}70 € je km ohne Grundgebühr.
\end{teile}
\begin{gleichungsraster}[2]
\gl{\text{Welches Mietauto? Angebot 1: 32 € Miete pro Tag, 300 km insgesamt frei, jeder weitere km 0{,}30 €. Angebot 2: einmalig 95 € für eine Woche, jeder km 0{,}12 €. Jana mietet 4 Tage und fährt 420 km. Berechne die Gesamtkosten beider Angebote, entscheide, welches günstiger ist, und stelle für Angebot 2 eine Gleichung für die Kosten y (in €) bei x km für eine Woche auf. (P10 2023 OS)}} & \\
\end{gleichungsraster}
\end{aufgabe}

%% TODO Titel: Ich-kann-Satz fehlt in der Bank (Platzhalter: Kettenname) – Nr. 29
%% TODO Anweisung: ein Satz über der ersten Teilaufgabe fehlt in der Bank – Nr. 29
\begin{aufgabe}{Graph zu Tarif zuordnen}
\begin{teile}
\teil Welcher Graph passt zum Tarif? Tarif: 10 € Grundgebühr und 2 € je Stunde. Gib den passenden Graphen an.

\begin{ksys}[xmin=0,xmax=8,ymin=0,ymax=40,xstep=1,ystep=5,ablesen] \gerade{2}{0}{A} \gerade{2}{10}{B} \gerade{10}{2}{C} \end{ksys}
Graph \leerfeld
\end{teile}
\end{aufgabe}

%% TODO Titel: Ich-kann-Satz fehlt in der Bank (Platzhalter: Kettenname) – Nr. 30
%% TODO Anweisung: ein Satz über der ersten Teilaufgabe fehlt in der Bank – Nr. 30
\begin{aufgabe}{Situation zu Gleichung beschreiben}
\begin{teile}
\teil Welche Situation passt? Beschreibe eine Situation zur Gleichung $y = 2{,}5x + 4$ und sage, was x und y bedeuten.
\end{teile}
\end{aufgabe}

%% TODO Titel: Ich-kann-Satz fehlt in der Bank (Platzhalter: Kettenname) – Nr. 31
%% TODO Anweisung: ein Satz über der ersten Teilaufgabe fehlt in der Bank – Nr. 31
\begin{aufgabe}{Anwendung – Fehler finden, Begründen, Darstellungswechsel, Anwendung}
\begin{teile}
\teil Ich finde den Fehler: Ein Verleih kostet 6 € Grundgebühr und 2{,}50 € je Stunde. Max schreibt $y = 6x + 2{,}5$. Benenne den Fehler und gib die richtige Gleichung an.
\teil Welcher Tarif ist immer günstiger? Tarif A: 10 € Grundgebühr und 0{,}20 € je Minute. Tarif B: 10 € Grundgebühr und 0{,}15 € je Minute. Begründe ohne Rechnung.
\teil Gleichung aus der Tabelle: Die Tabelle zeigt die Kosten y (in €) eines Handytarifs für x GB. Gib die Gleichung des Tarifs an.

\wertetabelle[12,15,18,21]{x}{y}{0,2,4,6}
y = \leerfeld
\end{teile}
\begin{gleichungsraster}[2]
\gl{\text{Wie weit reicht das Geld? Für die Klassenfahrt kostet ein Bus 180 € Grundpreis und 1{,}90 € je km. Die Klasse hat 600 € für den Bus. Wie viele Kilometer (hin und zurück) sind höchstens drin?}} & \\
\end{gleichungsraster}
\end{aufgabe}
